% Part: incompleteness
% Chapter: arithmetization-syntax
% Section: coding-formulas

\documentclass[../../../include/open-logic-section]{subfiles}

\begin{document}

\olfileid[id]{inc}{art}{frm}
\olsection{Pengodean \printtoken{P}{formula}}

Setelah kita mendefinisikan relasi $\fn{Term}(x)$ secara rekursif primitif,
kita dapat menggunakannya untuk mendefinisikan secara rekursif primitif relasi
yang bersesuaian bagi !!{formula}s, yakni $\fn{Frm}(x)$.

\begin{prop}
Relasi $\fn{Atom}(x)$, yang berlaku jika dan hanya jika $x$ adalah bilangan
G\"odel dari !!{formula} atomik, bersifat rekursif primitif.
\end{prop}

\begin{proof}
Bilangan $x$ adalah bilangan G\"odel dari !!{formula} atomik jika dan hanya jika
salah satu hal berikut berlaku:
\begin{enumerate}
\item Terdapat $n$, $j < x$, dan $z < x$ sedemikian rupa sehingga untuk setiap
  $i < n$, $\fn{Term}((z)_i)$, $\len{z}=n$, dan $x = $
\[
\Gn{\Obj P^n_j(} \concat \fn{flatten}(z) \concat \Gn{)}.
\]
\item Terdapat $z_1, z_2 < x$ sedemikian rupa sehingga $\fn{Term}(z_1)$,
  $\fn{Term}(z_2)$, dan $x = $
\[
\Gn{{\eq}(} \concat z_1 \concat \Gn{,} \concat z_2 \concat{\Gn{)}}.
\]
  \tagitem{prvFalse}{$x = \Gn{\lfalse}$.}{}
  \tagitem{prvTrue}{$x = \Gn{\ltrue}$.}{}
\end{enumerate}
\end{proof}

\begin{prop}
\ollabel{prop:frm-primrec}
Relasi $\fn{Frm}(x)$, yang berlaku jika dan hanya jika $x$ adalah bilangan
G\"odel !!a{formula}, bersifat rekursif primitif.
\end{prop}

\begin{proof}
Suatu barisan simbol~$s$ adalah !!a{formula} jika dan hanya jika terdapat
barisan pembentukan $s_0$, \dots, $s_{k-1} = s$ dari !!{formula} yang mencatat
bagaimana~$s$ dibentuk dari !!{formula}s atomik menurut aturan pembentukan.
Kode setiap $s_i$ tidak melebihi kode~$x$ dari~$s$. Selain itu, kita dapat
memilih $k \le \len{x}$, sehingga kode barisan
$\tuple{s_0, \dots, s_{k-1}}$ memiliki batas rekursif primitif,
misalnya $p_{k-1}^{k(x+1)}$. Karena itu, keberadaan barisan pembentukan yang
sesuai dapat dinyatakan dengan kuantifikasi terbatas rekursif primitif.
\end{proof}

\begin{prob}
Berikan pembuktian terperinci bagi
\olref[inc][art][frm]{prop:frm-primrec} dengan mengikuti pola pembuktian
pertama \olref[inc][art][trm]{prop:term-primrec}.
\end{prob}

\begin{prop}
\ollabel{prop:freeocc-primrec}
Relasi $\fn{FreeOcc}(x, z, i)$, yang berlaku jika dan hanya jika simbol ke-$i$
dalam formula dengan bilangan G\"odel~$x$ merupakan kemunculan bebas variabel
dengan bilangan G\"odel~$z$, bersifat rekursif primitif.
\end{prop}

\begin{proof}
Latihan.
\end{proof}

\begin{prob}
Buktikan \olref[inc][art][frm]{prop:freeocc-primrec}. Anda boleh menggunakan
fakta bahwa setiap subuntai dari !!a{formula} yang juga merupakan
!!a{formula} adalah sub-!!{formula} dari formula tersebut.
\end{prob}

\begin{prop}
Sifat $\fn{Sent}(x)$, yang berlaku jika dan hanya jika $x$ adalah bilangan
G\"odel !!a{sentence}, bersifat rekursif primitif.
\end{prop}

\begin{proof}
!!{sentence} adalah !!a{formula} tanpa kemunculan bebas !!{variable}s. Jadi,
$\fn{Sent}(x)$ berlaku jika dan hanya jika
\begin{multline*}
\fn{Frm}(x) \land {}\\
\bforall{i<\len{x}}{\bforall{z<x}{}}\\
(\bexists{j<z}{z=\Gn{\Obj v_j}} \lif \lnot\fn{FreeOcc}(x,z,i)).
\end{multline*}
\end{proof}


\end{document}
